% $Date: 2026/05/02 16:49:56 $
% This template file is public domain.
%
% TUGboat class documentation is at:
%   texdoc tugboat
% or
%   https://texdoc.org/serve/tugboat/0
% or
%   https://mirrors.ctan.org/macros/latex/contrib/tugboat/ltubguid.pdf

\documentclass{ltugboat}
\usepackage[T1]{fontenc}
\usepackage{graphicx}
\usepackage{microtype}
\usepackage[hidelinks]{hyperref}

\title{\LaTeX\ for Elegant Humanities Handouts in Conferences and Classrooms}

% repeat info for each author; comment out items that don't apply.
\author{Corey J. Stephan}
\address{2126 Berkhsire Elm Street \\ Katy, Texas, USA 77493 \\}
\netaddress{corey.stephan@stthom.edu}
\personalURL{https://coreystephan.com}
%\ORCID{0}

% Please state if you'd like to receive a physical copy of the TUGboat
% issue or if electronic access suffices.  If you want a physical issue,
% please include the mailing address we should use, as a comment if
% you prefer it not be printed.

% Physical copy, please, mailed here:
%
% Corey Stephan
% 2126 Berkshire Elm Street
% Katy, Texas, USA 77493
%

% Corey Stephan's package declarations
\usepackage{polyglossia}
\usepackage{parallel}
\usepackage{hologo}
\usepackage[olddefault]{fontsetup}


\begin{document}
	
% Corey Stephan's linguistic declarations
\setdefaultlanguage{english}
\setotherlanguage[variant=ancient]{greek}
	
\maketitle

\begin{abstract}
{\LaTeX\ has been popular throughout the mathematical and physical sciences for decades. \LaTeX's utility for scholars in humanities disciplines, however, tends to be overlooked. In this tutorial, I introduce the use of \LaTeX\ for a common task in humanities scholarship, namely, typesetting materials with unusual line length and other organizational needs, such as poetry and translations. The specific example is preparing a handout for a historical conference or class presentation that contains excerpts of ancient Greek with parallel English translations. Additionally, I provide general recommendations for learning and employing \LaTeX\ across the academic disciplines of the humanities. I aim to inspire colleagues in religious, historical, philosophical, and related areas of study to think of \LaTeX\ as having a wide range of potential professional applications even inside these scholarly fields in which the use of \LaTeX\ is not widespread.}
\end{abstract}

\section{\LaTeX\ for the humanities}

The world's best established software system for typesetting and document formatting, \LaTeX\ has been popular throughout the mathematical and physical sciences since before I was born. The utility of \LaTeX\ for humanities disciplines, however, long has been overlooked. The barrier to entry of learning a language for computer typesetting often means that persons outside disciplines in which \LaTeX\ is part of the academic culture are unlikely to bother. Plus, typically one begins to learn the main computing utilities that one will use in a given field (at least) as early as during one's undergruadate studies. If \LaTeX\ is not one of the computing utilities to which a student is exposed early in her/his formation, the student probably will not interact with \LaTeX\ later in her/his career. Yet, \LaTeX\ can be used for elegant typesetting in most instances in which a scholar of humanities needs to present materials, especially when the appearance of the materials matters per se. Whether one works in history, philosophy, literature, or some other discipline in which \LaTeX\ is not usually employed, \LaTeX\ is worth taking the time to learn.

Granted, in disciplines outside the mathematical and physical sciences, such as my own area of historical theology (technically, divinities rather than humanities), editors expect us to prepare our materials for publications with bug-for-bug compatibility with Microsoft Word so that we may have smooth back-and-forth markup. In practice, this expectation requires us to use word processors, such as LibreOffice Writer (among other free and open source options). A word processor presents material in a What You See Is What You Get (WYSIWYG, ``whizzy-wig'') style. While one writes, one simultaneously thinks about and tends to matters of formatting. (I suspect that this irksome question is familiar for everyone who publishes in the humanities: How do I need to adjust the spacing of this footnote so that it aligns with my University Press's Style Guide?)

In a typesetting language, such as \LaTeX\ (or \ConTeXt\ or Typst, either of which a scholar in humanities also might choose to use in various situations), one is more concerned about the structure of what one is writing than about the (final) formatting. One declares that specific textual items will be in certain categories of writing `things' (heading, paragraph, footnote, accent type, typeface, and so on), and one allows the program to compile one's writing `code' into a tidy final document \Dash (nearly) every time. Preparing a \textit{curriculum vitae}, conference handout, or other document with typesetting requirements beyond what are straightforward to achieve in a word processor becomes more about the writing than about the formatting. A significant investment of time might be required at the beginning of a project to configure a .tex file's preamble and/or other matters related to the `code' that will tell \LaTeX\ how to output a document from one's writing. Yet, once that investment has been made, one may simply write, declaring structural points as one proceeds.

To learn \LaTeX\ as a scholar outside the disciplines in which \LaTeX\ is the norm, I recommend working through one of two community-built, libre guidebooks. The first, which I used to learn \LaTeX, is \titleref{The Not So Short Introduction to \LaTeX} (\url{github.com/oetiker/lshort}). The second, which a colleague of mine recommends, is \titleref{Modern \LaTeX} (\url{github.com/mrkline/modern-latex}). Either guidebook may be used to begin to learn to typeset cover letters, ancient languages, and other materials that are relevant across the humanities \Dash without the front-loading of mathematics and physical sciences that one finds in most other introductions. Moreover, during the Q\&A at the end of the talk at \TUG\ 2026 on which this article is based, Takuto Asakura helpfully added \url{learnlatex.org}, which is an official part of the \LaTeX\ Project, as a third recommendation. This website provides an interactive tutorial that anyone may use who intends to start using \LaTeX\ for any purpose. All of its learning materials are available in several major world languages (including Asakura's native Japanese).

To show how thinking about digital writing in terms of document structure rather than final format can be of practical use in the humanities, here I present a tutorial in which I introduce the use of \LaTeX\ for a common task in humanities scholarship, namely, typesetting texts with unusual line length and other organizational needs. The case study is preparing a handout for a historical conference or class presentation with perfectly organized excerpts of a Byzantine Greek text and an original English translation in parallel. The principles that I outline will be applicable to any presentation of parallel text and, more generally, most varieties of handouts that are popular across the humanities.

\section{Simplest side-by-side text in \LaTeX: \texttt{parallel}}

There are many ways to set parallel text in \LaTeX. The simplest is to use the community package called \texttt{parallel}, which is included in \hologo{TeXLive}, \hologo{MiKTeX}, and other \LaTeX\ distrubtions. For full documentation, visit \url{ctan.org/pkg/parallel}.

To start using \texttt{parallel} in any \LaTeX\ document, include the following in the preamble:

\begin{verbatim}
\usepackage{parallel}
\end{verbatim}
For example, here is an abbreviation of the preamble to the Beamer (\LaTeX\ slide show) file for my talk at the 2026 conference of the \TUG\ (\url{github.com/historical-theology/latex-talks/}) on which this article is based with only the most important entries, including the package \texttt{parallel}:

\begin{verbatim}
\documentclass[aspectratio=169]{beamer}
\usetheme{Goettingen}
...
\usepackage{parallel}
...
\title{\LaTeX\ for Elegant 
	Humanities Handouts \\ 
	in Conferences and Classrooms }
\author{Corey Stephan, Ph.D.}
...
\begin{document}
\end{verbatim}
I begin by declaring that the class of the document is \texttt{beamer} (with a 16:9 aspect ratio for all slides, per the specifications of the \TUG). Then, after some more (suppressed) entries, I actually declare that I will be using the package \texttt{parallel}. After further (suppressed) entries, I assign a title to the document and supply my name in the designated place for the name of the author. Finally, I trigger the document (writing) portion to begin.

This is the basic format for parallel text set with the package \texttt{parallel}:

\begin{verbatim}
\begin{Parallel}{}<L space>}{<R space>}
\ParallelLText{<left text row 1>}
\ParallelRText{<right text row 1>}
\ParallelPar
\ParallelLText{<left text row 2>}
[and so on]
\end{Parallel}	
\end{verbatim}
First, one opens an instance of \texttt{parallel}. Then, one writes the first left row, followed by the adjacent right row. One triggers a new line with \verb|\ParallelPar\]| (an abbreviation of \textit{parallel paragraph}). Then, one writes the second left row, followed by the second right row, and so on until one ends the instance of \texttt{parallel}.

I (a theologian) might typeset the opening words of the Christian prayer \textit{Ave Maria} in their traditional Latin with parallel English like so:

\begin{verbatim}
\begin{Parallel}{0em}{0em}
	\ParallelLText{Ave Maria,}
	\ParallelRText{Hail Mary,}
		\ParallelPar
	\ParallelLText{gratia plena \ldots}
	\ParallelRText{full of grace \ldots}
		\ParallelPar
\end{Parallel}	
\end{verbatim}

Here is the result:

\medskip
\begin{minipage}{.4\textwidth}
\begin{Parallel}{0em}{0em}
\ParallelLText{Ave Maria,}
\ParallelRText{Hail Mary,}
\ParallelPar
\ParallelLText{gratia plena \ldots}
\ParallelRText{full of grace \ldots}
\ParallelPar
\end{Parallel}
\end{minipage}\\

\texttt{parallel} makes it easy to set short sections of text, from a few lines to a few stanzas or paragraphs, side-by-side in \LaTeX.

\section{Advanced parallel text in \LaTeX:\\ \texttt{reledpar} and \texttt{reledmac}}

For more advanced setting of parallel text in \LaTeX, especially long-form (e.g. a critical edition with accompanying translation), one ought to learn the two interconnected community packages \texttt{reledpar} (for parallel typesetting specifically) and \texttt{reledmac} (for typesetting academic editions in general), which are bundled with the major \LaTeX\ distributions.

\texttt{reledpar} and \texttt{reledmac} are purpose-built for long-form academic typesetting, especially in academic articles, books, and critical editions. These two packages are meant to replace older methods of preparing academic editions and parallel texts in \LaTeX\ that one might find in Web searches (notably, the outdated \texttt{ledpar} and \texttt{ledmac} packages, which \texttt{reledpar} and \texttt{reledmac} replace).

\texttt{reledpar} and \texttt{reledmac} also offer flexible functionality that is difficult to find in any software outside \LaTeX, such as \texttt{reledmac}'s ability to wrangle the apparatus of a textual critical edition with mere markup syntax. (At \TUG\ 2026, Ulrike Fischer informed me that she has interacted with German-speaking scholars who use \LaTeX\ precisely for its smooth flexibility with typesetting critical editions.) For the full documentation of these two packages, visit \url{ctan.org/pkg/reledpar} and \url{ctan.org/pkg/reledmac}.

\section{Setting non-Latin characters in \LaTeX: \texttt{babel} and \texttt{polyglossia}}

The most important \LaTeX\ packages for languages with non-Latin alphabets are \texttt{babel} and \texttt{polyglossia}, which are included in the major \LaTeX\ distributions. Although it is possible to convert a project from one of these metapackages to the other (since the syntax is highly similar between the two), a scholar ought to invest the time to choose carefully between either of them at the start of a grand project.

In \titleref{The \LaTeX\ Companion: 3rd Edition} (Addison-Wesley, 2023), Frank Mittelbach and Ulrike Fischer suggest that despite the fact that \texttt{babel} is older than \texttt{polyglossia}, \texttt{babel} generally is preferrable: ``Unless you are in a special situation where a language or script is notably better supported by \texttt{polyglossia}, the recommendation is to use \texttt{babel} with all engines'' (entry 13.7.2, ``Polyglossia''). Since \texttt{babel} technically supports the same languages as \texttt{polyglossia}, a situation in which one would need \texttt{polyglossia} over \texttt{babel} is rare. Yet, a scholar working with an ancient language might face a situation in which, say, \texttt{polyglossia} actually supports that language better than \texttt{babel}. To make an informed choice against a given academic undertaking's particular requirements, one should explore each package's documentation at \url{ctan.org/pkg/babel} and \url{ctan.org/pkg/polyglossia} \Dash and perhaps even run a small test with both \texttt{babel} and \texttt{polyglossia} with samples of the non-Latin writing that one intends to typeset.

In (old) Web forum discussions, one might read that \texttt{babel} will work better with the \LaTeX\ engine \LuaLaTeX, whereas \texttt{polyglossia} will work better with \XeLaTeX. As of the time of writing this essay, however, one may choose either of the two mainline \LaTeX\ engines that are built with full Unicode support, \XeLaTeX\ or \LuaLaTeX, for either \texttt{babel} or \texttt{polyglossia}. Results might vary depending on the specifics of a given project.

To learn \texttt{babel}, I recommend Chapter 13: ``Localizing documents'' in \titleref{The \LaTeX\ Companion: 3rd Edition}. Here, Mittelbach and Fischer succinctly introduce localization in \LaTeX\ via \texttt{babel} in a practical, didactic style. Of course, the package's official documentation in \CTAN\ also ought to be consulted for assistance with complex arrangements.

To learn \texttt{polyglossia}, I recommend the package's documentation in \CTAN.

\section{Starting a project with \texttt{polyglossia}}

Here, I will model \texttt{polyglossia}, which is what I happen to be familiar with using for polytonic (ancient and Byzantine) Greek.

First, include the package in the preamble:
\begin{verbatim}
\usepackage{polyglossia}
\end{verbatim} 
Then, after the beginning of the document but before actual type, declare the languages and their variants. For example, here are the \texttt{polyglossia} declarations for the .tex file that was compiled into this document:

\begin{verbatim}
\begin{document}
	\setdefaultlanguage{english}
	\setotherlanguage[variant=ancient]{greek}
\end{verbatim}
It is important to explore the \texttt{variant} options that are available for a given language in the official \texttt{polyglossia} (and \texttt{babel}) documentation, since setting this to the right value enables correct hyphenation, accent marks, and other matters of typography. Here, I use \texttt{variant=ancient} to set polytonic Greek rather than the default Greek variant (i.e. Modern).

One is able to use one of the `other' declared language(s) (in this case, ancient Greek) by toggling it when needed.

A declared foreign or ancient language may be toggled for a long excerpt like so:

\begin{verbatim}
\begin{greek}
	<greek text>
	...
\end{greek}
\end{verbatim}
For example, this is how one could set the famous opening verse of John's Gospel:

\begin{verbatim}
\begin{greek}
ἐν ἀρχῇ ἦν ὁ λόγος, \\
καὶ ὁ λόγος ἦν πρὸς τὸν θεόν, \\
καὶ θεὸς ἦν ὁ λόγος. \\
\end{greek}
\end{verbatim}
%
Here is the clean result:
%
\begin{greek} \\
ἐν ἀρχῇ ἦν ὁ λόγος, \\
καὶ ὁ λόγος ἦν πρὸς τὸν θεόν, \\
καὶ θεὸς ἦν ὁ λόγος. \\
\end{greek} 
\\
For another example, these are the first two lines of Homer's \titleref{The Iliad}:

\begin{verbatim}
\begin{greek}
μῆνιν ἄειδε θεὰ Πηληϊάδεω Ἀχιλῆος \\
οὐλομένην, ἣ μυρί' Ἀχαιοῖς ἄλγε' ἔθηκε,
\end{greek}
\end{verbatim}
%
The result is equally clean to that of John 1:1, even for these particular words of Homer with their more complicated accent and breathing marks:
%
\begin{greek}\\
μῆνιν ἄειδε θεὰ Πηληϊάδεω Ἀχιλῆος \\
οὐλομένην, ἣ μυρί' Ἀχαιοῖς ἄλγε' ἔθηκε,
\end{greek} 
\\


A declared language may be toggled for a single word or phrase (rather than a series of verses or lines) in this way:

\begin{verbatim}
\text<language>{<text in language>}
\end{verbatim}
%
This mid-line toggling style may be used to insert a keyword or short quotation in that declared language, e.g. {\textgreek{ὁ λόγος}} (produced with \verb|\textgreek{ὁ λόγος}|), inside running prose in another language, such as the document's primary one.

\section{Combining \texttt{parallel} and \texttt{polyglossia}}

The simplest way of setting multilingual text side-by-side is with the package \texttt{parallel}, as I showed in section (2). Since \texttt{parallel} can be tedious to organize for more than a few stanzas/paragraphs, I remind the reader of the packages \texttt{reledpar} and \texttt{reledmac} for long and/or complicated parallel text displays, as I suggested in section (3).

Here is an example of \texttt{parallel} and \texttt{polyglossia} working in combination to display a key excerpt of Greek text with my own English translation from the opening chapter of my recently published historical monograph \titleref{Maximus the Confessor's Thomistic Legacy}. The idea is that I might present the material in this way on a conference or classroom handout:

\begin{verbatim}
\begin{Parallel}{0em}{0em}
\textgreek{\ParallelLText
	{Νῦν συγχορεύεις ἀληθῶς}}
\ParallelRText
	{Now, truly, you celebrate together}
	\ParallelPar
[and so on]
\end{Parallel}	
\end{verbatim}
The result might look like Figure 1: ``Canon in Honor of Thomas Aquinas,'' Ode 3, which spans the full width of the next page.\footnote{The Greek text in Figure 1 is from Raffaele Cantarella, ed., “Canone greco inedito di Giuseppe vescovo di Methone in onore di San Tommaso d’Aquino,” \titleref{Archivum Fratrum Pr\ae dicatorum} 12 (1934): 145--85, at 155--58. This parallel Greek text with my Latin translation duplicates Corey J. Stephan, \titleref{Maximus the Confessor’s Thomistic Legacy}, Medi\ae val Law and Theology 10 / Studies and Texts 244 (Toronto: Pontifical Institute of Medi\ae val Studies, 2026), 6.} With \texttt{parallel}, all spacings \Dash between lines, letters, and otherwise \Dash appear exactly how one expects without any of the (dreadful) fiddling that comes with trying to achieve a similar result using tables in a word processor.

\begin{figure*}
\begin{Parallel}{0em}{0em}
\textgreek{\ParallelLText{Νῦν συγχορεύεις ἀληθῶς}}
\ParallelRText{Now, truly, you celebrate together}
\ParallelPar
\textgreek{\ParallelLText{σὺν τοῖς λοιποῖς θεολόγοις}}
\ParallelRText{with the rest of the theologians}
\ParallelPar
\textgreek{\ParallelLText{καὶ τοῖς ἄλλοις θεολήπτοις ἀνδράσι,}}
\ParallelRText{and the other divinely inspired men}
\ParallelPar
\ParallelRText{with the rest of the theologians}
\textgreek{\ParallelLText{περὶ τὸν θρόνον Χριστοῦ}}
\ParallelRText{around the throne of Christ,}
\ParallelPar
\textgreek{\ParallelLText{τὸν ἱερὸν ἐξᾴδοντες}}
\ParallelRText{singing the sacred}
\ParallelPar
\textgreek{\ParallelLText{καὶ τρισάγιον ὕμνον}}
\ParallelRText{and Thrice-Holy Hymn}
\ParallelPar
\textgreek{\ParallelLText{σὺν ταῖς ἁγίαις δυνάμεσιν.}}
\ParallelRText{with the holy powers.}
\ParallelPar

\vspace{1em}

\textgreek{\ParallelLText{Ἔνθα πατέρων ὁ χορὸς}}
\ParallelRText{Where there is the chorus of Fathers}
\ParallelPar
\textgreek{\ParallelLText{καὶ ἱερῶν θεολόγων}}
\ParallelRText{and holy theologians,}
\ParallelPar
\textgreek{\ParallelLText{ὡς ὑψίνους καὶ λεπτὸς θεόλογος}}
\ParallelRText{as a high-minded and subtle theologian}
\ParallelPar
\textgreek{\ParallelLText{κατεσκήνωσας, Θωμᾶ,}}
\ParallelRText{you do abide, Thomas,}
\ParallelPar
\textgreek{\ParallelLText{σὺν τούτοις αὐλιζόμενος,}}
\ParallelRText{dwelling with them,}
\ParallelPar
\textgreek{\ParallelLText{θεολογῶν ἀπαύστως}}
\ParallelRText{unceasingly theologizing,}
\ParallelPar
\textgreek{\ParallelLText{τριάδα, πάτερ, τὴν ἄκτιστον.}}
\ParallelRText{O Father, about the Uncreated Trinity.}
\ParallelPar

\vspace{1em}

\textgreek{\ParallelLText{Πῦρ ἀνεδείχθης φλογερὸν}}
\ParallelRText{A burning fire you have shown yourself to be}
\ParallelPar
\textgreek{\ParallelLText{τὰς φρυγανώδεις αἱρέσεις}}
\ParallelRText{against the firewood-like heresies;}
\ParallelPar
\textgreek{\ParallelLText{τῇ ἱσχύϊ τῆς σοφίας σου, πάτερ,}}
\ParallelRText{by the might of your wisdom, O Father,}
\ParallelPar
\textgreek{\ParallelLText{κατακαύσας, καὶ πιστῶν}}
\ParallelRText{you have burned them down, and}
\ParallelPar
\textgreek{\ParallelLText{καταφαιδρύνας, ἅγιε,}}
\ParallelRText{you have restored completely, O Saint,}
\ParallelPar
\textgreek{\ParallelLText{τὰς διανοίας ὅθεν}}
\ParallelRText{the thoughts of the faithful, whence}
\ParallelPar
\textgreek{\ParallelLText{σὲ συνελθόντες γεραίρομεν.}}
\ParallelRText{we, all joined together, do venerate you.}
\ParallelPar
\end{Parallel}
\caption{``Canon in Honor of Thomas Aquinas,'' Ode 3}
\end{figure*}

\section{Other \LaTeX\ tools for humanities handouts}

There are many \LaTeX\ tools that can assist work across the humanities disciplines. Here is a succinct, curated list of the most broadly applicable, especially those that one might wish to use to assist with preparing a handout for a conference presentation or university class period:

\begin{itemize}
	\item Either \BibLaTeX\ or the older but still useful \BibTeX\ offers a comprehensive citation and bibliography management system that is capable of being interlinked with a variety of citation managers \Dash notably, for academic research, Zotero and JabRef.
		\subitem \BibLaTeX\ and \BibTeX\ are both built into \hologo{TeXLive} and other \LaTeX\ distributions.
	
	\item The free and open source (GPL-2) graphical software application \texttt{TeXstudio} (\url{texstudio.org}) offers a \LaTeX\ -specific writing environment with comfortable features like intuitive autocompletion of \LaTeX\ commands, many options for custom theming, and even a (self-hosted) collaborative (online) editing option.
		\subitem \LaTeX\ may be written in any plain text editor, but a designated environment improves efficiency and helps avoid code errors.
		\subitem I specifically mention \texttt{TeXstudio} for two reasons. First, it is what I use for my own work. Second, on a practical note, it is included in the official software repositories of most libre Unix-like operating systems (GNU/Linux and BSD), making it easy to install in a context of software freedom.
	
	\item By running search queries in \CTAN, one may find community-built packages for many individual languages that are of relevance for historical or other academic research.
		\subitem An example is \texttt{aramaic-serto} for the ancient Middle-Eastern language Syriac (Aramaic).
	
	\item \texttt{microtype} offers many small adjustments to type, such as improved automatic justification and hyphenation.
		\subitem I use \texttt{microtype} in everything that I make professionally in \LaTeX. Atop \LaTeX's default elegance, \texttt{microtype} helps self-produced documents look as similar to materials that have been prepared by a press's typesetter as is reasonably possible.
		\subitem This TUGboat article provides a real-world example of the \texttt{microtype} package in-use.
\end{itemize}





\makesignature
\end{document}
